\capitulo{CAPÍTULO 1}{La responsabilidad extracontractual y su diferencia con otros regímenes de responsabilidad}

\section{1.1. ¿Qué es la responsabilidad?}
Un texto ilustra los distintos sentidos que tiene la palabra «responsabilidad» en el lenguaje común:

\begin{cita}
\textit{«Como capitán de un barco, X era responsable por la seguridad de sus pasajeros y su carga. Pero en su último viaje, se embriagaba todas las noches y fue responsable de la pérdida del barco con todo lo que llevaba. Se rumoraba que estaba loco, pero los médicos lo encontraron responsable de todas sus acciones. Durante el viaje, X se comportó muy irresponsablemente y varios incidentes que tuvo en su carrera demuestran que no era una persona responsable. El capitán siempre sostuvo que fueron las tormentas excepcionales las responsables por la pérdida del barco, pero en un proceso judicial que se le siguió fue encontrado responsable por la pérdida de vidas y bienes. Todavía vive y es moralmente responsable por la muerte de muchas mujeres y niños.»}
\end{cita}
\begin{ejemplo}
\textcolor[HTML]{4F6228}{\textbf{Ejemplo. }}En síntesis: el capitán es «responsable» por la seguridad de sus pasajeros; se embriaga todas las noches y es «responsable» de la pérdida del barco; los médicos lo encuentran «responsable» de sus actos; se comporta «irresponsablemente»; él culpa a las tormentas como «responsables» de la pérdida; un juez lo declara «responsable» de las vidas y bienes perdidos, y sigue siendo «moralmente responsable» de las muertes.
\end{ejemplo}
Sentidos de la palabra «responsabilidad» que muestra el ejemplo:

\begin{itemize}
\item Obligaciones o deberes derivados de un cargo.
\item Factor causal.
\item Estado mental e imputabilidad.
\item Estándar de diligencia.
\item Calificación de conductas reprochables.
\end{itemize}
\section{1.2. Evolución histórica de la responsabilidad extracontractual}
\subsection{1.2.1. Derecho romano: la Ley de las XII Tablas}
\begin{cita}
\textit{«Tab. 8, 10. El que hubiere incendiado una casa o un montón de trigo puesto junto a la casa, se manda que atado y azotado sea quemado vivo, si a sabiendas y deliberadamente hubiere hecho esto; pero si por casualidad, esto es, por negligencia, se manda o que resarza el daño, o, si no fuera solvente, es castigado con más lenidad; más en la denominación de la casa se comprenden todas las especies de edificios.»}
\par
\textit{«Tab. 8, 2. Si se ha amputado un miembro y no se logra la composición pecuniaria, se aplique el talión.»}
\par
\textit{«Tab. 8, 3. Si se ha fracturado un hueso a un hombre libre, está obligado a pagar la pena de 300 sestercios, si es un esclavo, de 150.»}
\par
\textit{«Tab. 8, 4. Si se injuria a otro, la pena es de 25 sestercios.»}
\end{cita}
En otras palabras: quien incendia una casa o el trigo junto a ella a sabiendas es azotado y quemado vivo; si lo hace por negligencia, debe resarcir el daño o, si es insolvente, recibe un castigo más leve. Si se amputa un miembro y no hay composición pecuniaria, se aplica el talión. Fracturar un hueso a un hombre libre se paga con 300 sestercios; a un esclavo, con 150. La injuria se paga con 25 sestercios.

Características:

\begin{itemize}
\item Distingue entre daños a las cosas y daños a las personas.
\item Es un derecho casuístico.
\item Confunde la responsabilidad penal con la extracontractual (confusión entre pena y reparación).
\item La \textit{iniuria} se entiende como acción (delito) o conducta antijurídica.
\end{itemize}
\subsection{1.2.2. Derecho romano: la \textit{Lex Aquilia}}
\begin{cita}
\textit{«El que hubiere matado con injuria al esclavo o a la esclava ajenos, a un cuadrúpedo, o a una res, sea condenado a pagar al dueño el precio mayor que aquello tuvo en aquel año.»}
\par
\textit{«Respecto a las demás cosas, excepto el esclavo y las reses que hayan sido muertos, si alguien hiciere daño a otro, porque hubiere quemado, quebrado o roto alguna cosa con injuria, sea condenado a pagar al dueño tanto cuanto aquella cosa valiere en los treinta días próximos.»}
\end{cita}
Es decir: quien mata con injuria a un esclavo o a un animal ajeno paga al dueño el mayor valor que tuvo en ese año; por las demás cosas que se quemen, quiebren o rompan con injuria, se paga lo que valían en los treinta días siguientes.

Características:

\begin{itemize}
\item La \textit{iniuria} pasa a ser una calificación de la conducta: primero el dolo, luego la culpa.
\item Surge una cláusula general de daño.
\item Surgen criterios de valoración de perjuicios.
\item Persisten la \textit{iniuria} como delito y la confusión entre responsabilidad penal y extracontractual (entre pena y reparación).
\end{itemize}
\subsection{1.2.3. Derecho romano: la acción edictal (\textit{actio de effusis vel deiectis})}
\begin{cita}
\textit{«Si se hubiere arrojado o derramado alguna cosa en el sitio por donde vulgarmente se transita, o donde la gente se detiene, daré, contra el que allí habitare, acción en el duplo por cuanto daño con ello se hubiere causado o hecho. Si se dijera que del golpe de lo arrojado había perecido un hombre libre, daré acción de cincuenta áureos; si viviera, y se dijese que se le causó daño, daré acción para aquel contra quien se reclama sea condenado en tanto cuanto por tal cosa pareciera justo al juez.»}
\end{cita}
Es decir: si se arroja o derrama algo en un lugar de tránsito, se concede acción contra quien habita allí por el doble del daño; si muere un hombre libre, la acción es de cincuenta áureos; si sobrevive, por lo que el juez considere justo.

Características:

\begin{itemize}
\item Surgen las primeras hipótesis de responsabilidad objetiva.
\item Nace del crecimiento urbanístico de Roma.
\item Responde a la necesidad de indemnizar cuando es difícil individualizar al responsable.
\item Se funda en la condición de \textit{habitator} (quien habita el lugar).
\item Sigue la lógica económica del principio \textit{cuius commoda eius et incommoda}: quien se beneficia de algo también debe asumir sus inconvenientes.
\end{itemize}
\subsection{1.2.4. Período moderno (anterior al Código Civil): el iusnaturalismo}
En el período moderno, el iusnaturalismo, con autores como Grocio, contribuyó al abandono progresivo del derecho casuístico y a la expansión de la idea de daño. La culpa empezó a consolidarse como fundamento de la responsabilidad.

\begin{cita}
\textit{«Llamamos aquí }maleficium\textit{ a toda culpa, ya en el hacer, ya en no hacer, que pugna con aquello que los hombres deben hacer comúnmente o por razón de cierta cualidad. De tal culpa nace, naturalmente, obligación, si se hizo daño, de que se resarza.» (Grocio)}
\end{cita}
Características del iusnaturalismo:

\begin{itemize}
\item La culpa como fundamento de la responsabilidad.
\item Abandono del derecho casuístico y expansión del daño.
\end{itemize}
Características del período premoderno (anterior a los códigos civiles):

\begin{itemize}
\item Simbiosis entre la moral cristiana y el derecho romano.
\item La culpa como fundamento de la responsabilidad.
\item Diferenciación entre la responsabilidad civil y la penal.
\end{itemize}
\subsection{1.2.5. Surgimiento de los códigos civiles}
Con la codificación quedan diferenciadas tres clases de responsabilidad:

\begin{itemize}
\item Responsabilidad penal.
\item Responsabilidad contractual.
\item Responsabilidad extracontractual.
\end{itemize}
\section{1.3. Responsabilidad subjetiva y responsabilidad objetiva}
\begin{definicion}
\textcolor[HTML]{1F3864}{\textbf{Definición. }}\textbf{Responsabilidad subjetiva: }aquella en la que, para atribuir responsabilidad, debe probarse que el autor del daño actuó con culpa o dolo, además de probarse el daño y el nexo causal.
\end{definicion}
\begin{formula}
{\large \textbf{Régimen subjetivo = daño + nexo causal + culpa}}
\end{formula}
\begin{definicion}
\textcolor[HTML]{1F3864}{\textbf{Definición. }}\textbf{Responsabilidad objetiva: }no exige probar culpa o dolo; basta acreditar el daño y su imputación conforme al régimen jurídico aplicable.
\end{definicion}
\begin{formula}
{\large \textbf{Régimen objetivo = daño + nexo causal (o relación de imputación)}}
\end{formula}
En el régimen subjetivo deben estar los tres elementos para que haya condena; en el objetivo, dos. En Colombia, la falla del servicio responde a una lógica subjetiva, mientras que títulos como el riesgo excepcional y el daño especial corresponden a esquemas objetivos.

\begin{comentario}
\textcolor[HTML]{8A6D3B}{\textbf{Comentario. }}Debe distinguirse entre los elementos que estructuran la responsabilidad (tres en el régimen subjetivo, dos en el objetivo) y quién tiene la carga de probarlos: son problemas diferentes.
\end{comentario}
\section{1.4. Concepto de responsabilidad extracontractual}
\begin{definicion}
\textcolor[HTML]{1F3864}{\textbf{Definición. }}\textbf{Responsabilidad extracontractual: }es la obligación de reparar daños causados por hechos jurídicos. Es una fuente de obligaciones en materia civil y surge por daños producidos por hechos distintos de los contratos.
\end{definicion}
\textbf{Diferencia entre acto y hecho jurídico:}

\begin{definicion}
\textcolor[HTML]{1F3864}{\textbf{Definición. }}\textbf{Acto jurídico: }manifestación de voluntad encaminada a producir efectos jurídicos, o acuerdo de voluntades entre sujetos que buscan generar efectos jurídicos.
\par
\textcolor[HTML]{1F3864}{\textbf{Definición. }}\textbf{Hecho jurídico: }fenómeno humano o no humano que produce consecuencias jurídicas sin que exista o medie necesariamente la voluntad de producir ese efecto jurídico concreto.
\end{definicion}
\section{1.5. Responsabilidad extracontractual y responsabilidad contractual}
\subsection{1.5.1. Diferencias}
\begin{tabla}
\begin{longtable}{|>{\raggedright\arraybackslash}m{\dimexpr 0.152\linewidth-2\tabcolsep-0.533pt\relax}|>{\raggedright\arraybackslash}m{\dimexpr 0.364\linewidth-2\tabcolsep-0.533pt\relax}|>{\raggedright\arraybackslash}m{\dimexpr 0.484\linewidth-2\tabcolsep-0.533pt\relax}|}
\hline
\cellcolor[HTML]{1F3864}\textcolor[HTML]{FFFFFF}{\textbf{Criterio}} & \cellcolor[HTML]{1F3864}\textcolor[HTML]{FFFFFF}{\textbf{Responsabilidad extracontractual}} & \cellcolor[HTML]{1F3864}\textcolor[HTML]{FFFFFF}{\textbf{Responsabilidad contractual}} \\ \hline
Fuente del daño & Hecho jurídico & Acto jurídico o contrato \\ \hline
\cellcolor[HTML]{F2F5FA}Elementos & \cellcolor[HTML]{F2F5FA}Según el régimen: subjetivo (daño, nexo y culpa) u objetivo (daño y nexo causal o imputación) & \cellcolor[HTML]{F2F5FA}Contrato + incumplimiento + daño + nexo causal + imputación \\ \hline
Indemnización & Reparación integral del daño & Por regla general, solo los daños previsibles; si hay dolo, también los perjuicios inmediatos o directos (art. 1616 del C.C.) \\ \hline
\end{longtable}
\end{tabla}
\textbf{Elementos en la responsabilidad contractual.} Sí se exige imputación al deudor; la culpa puede integrar ese juicio de imputación, pero no siempre debe ser probada por el acreedor, pues en ciertos eventos se presume con el incumplimiento.

\begin{itemize}
\item \textbf{Obligaciones de medio:} la responsabilidad contractual suele ser subjetiva; se analiza si el deudor actuó con la diligencia debida.
\item \textbf{Obligaciones de resultado:} la responsabilidad contractual es objetiva; no hay que demostrar si se incumplió con o sin culpa: basta, en principio, demostrar que el resultado prometido no se obtuvo, y el deudor debe acreditar una causa que lo exonere.
\end{itemize}
\textbf{Indemnización en la responsabilidad contractual.} Solo se reclaman los perjuicios previsibles, a menos que haya dolo (art. 1616 del C.C.).

\subsection{1.5.2. Similitudes}
En ambos ámbitos puede existir responsabilidad subjetiva u objetiva, según la obligación y el régimen aplicable:

\begin{itemize}
\item Régimen subjetivo de responsabilidad.
\item Régimen objetivo de responsabilidad.
\end{itemize}
\section{1.6. Responsabilidad extracontractual y responsabilidad penal}
\subsection{1.6.1. Diferencias}
\begin{tabla}
\begin{longtable}{|>{\raggedright\arraybackslash}m{\dimexpr 0.142\linewidth-2\tabcolsep-0.533pt\relax}|>{\raggedright\arraybackslash}m{\dimexpr 0.454\linewidth-2\tabcolsep-0.533pt\relax}|>{\raggedright\arraybackslash}m{\dimexpr 0.404\linewidth-2\tabcolsep-0.533pt\relax}|}
\hline
\cellcolor[HTML]{1F3864}\textcolor[HTML]{FFFFFF}{\textbf{Criterio}} & \cellcolor[HTML]{1F3864}\textcolor[HTML]{FFFFFF}{\textbf{Responsabilidad extracontractual}} & \cellcolor[HTML]{1F3864}\textcolor[HTML]{FFFFFF}{\textbf{Responsabilidad penal}} \\ \hline
Elementos & Subjetivo: daño, nexo y culpa. Objetivo: daño y nexo causal o imputación & Tipicidad, antijuridicidad y culpabilidad \\ \hline
\cellcolor[HTML]{F2F5FA}Función (¿para qué existe?) & \cellcolor[HTML]{F2F5FA}Principalmente resarcitoria o indemnizatoria: reparar el daño injustamente causado y restablecer, en lo posible, el equilibrio alterado & \cellcolor[HTML]{F2F5FA}Retributiva, sancionatoria, preventiva general y especial, y resocializadora \\ \hline
Consecuencia & Condena indemnizatoria (monetaria) & Penas previamente tipificadas en el Código Penal, incluida la privación o restricción de la libertad, y penas accesorias \\ \hline
\end{longtable}
\end{tabla}
\textbf{La función de la responsabilidad extracontractual.} Busca reparar el daño injustamente causado a una persona y restablecer, en la medida de lo posible, el equilibrio alterado por ese daño. Su función esencial es resarcitoria o indemnizatoria: que la víctima sea compensada por los perjuicios sufridos. Además, se le atribuye una función preventiva (incentiva a las personas y a las autoridades a actuar con diligencia para evitar daños) y una función de distribución de riesgos, especialmente en regímenes de responsabilidad objetiva. En el caso del Estado colombiano, esta función se conecta con el artículo 90 de la Constitución, que impone la obligación de reparar los daños antijurídicos que le sean imputables. (Sobre las objeciones a la función preventiva, ver Capítulo 2, sección 2.8.)

\begin{precision}
\textcolor[HTML]{9C2B2B}{\textbf{Precisión. }}La consecuencia de la responsabilidad extracontractual es una \textbf{condena indemnizatoria (resarcitoria)}, no sancionatoria. En Colombia la responsabilidad civil extracontractual y la del Estado buscan reparar el daño, no castigar: el artículo 16 de la Ley 446 de 1998 consagra el principio de reparación integral («todo el daño, nada más que el daño») y nuestro ordenamiento no reconoce, como regla general, los daños punitivos. La medida de la condena es el perjuicio probado, no la gravedad de la conducta. Lo «sancionatorio» es propio de la pena (responsabilidad penal). Las excepciones discutibles son la posibilidad de triplicar los perjuicios morales en graves violaciones de derechos humanos (Capítulo 4) y la multa de hasta 1.000 SMLMV de la Ley 2195 de 2022 (Capítulo 9), que introducen una función sancionatoria.
\end{precision}
\subsection{1.6.2. Similitudes}
\begin{itemize}
\item \textbf{La fuente de la responsabilidad:} un mismo evento puede activar los dos mundos. El delito es, al mismo tiempo, un hecho jurídico generador de responsabilidad civil extracontractual; por eso se activan ambas responsabilidades, la penal y la civil extracontractual.
\end{itemize}
\begin{ejemplo}
\textcolor[HTML]{4F6228}{\textbf{Ejemplo. }}El caso O. J. Simpson: fue absuelto en el proceso penal (por duda), pero perdió el proceso civil y fue condenado a pagar por el homicidio de su esposa.
\end{ejemplo}
\begin{itemize}
\item \textbf{El estudio de la causalidad:} en ambos campos es relevante.
\end{itemize}
\begin{comentario}
\textcolor[HTML]{8A6D3B}{\textbf{Comentario. }}Las teorías sobre la causalidad (equivalencia de las condiciones, causa próxima, causalidad adecuada e imputación objetiva) se originaron en el derecho penal y luego fueron trasladadas por los civilistas y por la jurisdicción de lo contencioso administrativo a la responsabilidad extracontractual (ver Capítulo 5). Además, la acción penal y la acción de reparación directa son acciones diferentes (ver Capítulo 4).
\end{comentario}
